\documentclass[conference]{IEEEtran}
\IEEEoverridecommandlockouts
\usepackage{cite}
\usepackage{amsmath,amssymb,amsfonts}
\usepackage{algorithmic}
\usepackage{graphicx}
\usepackage{textcomp}
\usepackage{xcolor}
\usepackage{booktabs}
\usepackage{hyperref}
\def\BibTeX{{\rm B\kern-.05em{\sc i\kern-.025em b}\kern-.08em
    T\kern-.1667em\lower.7ex\hbox{E}\kern-.125emX}}

\begin{document}

\title{WPipe: A Declarative Pipeline Framework with Forensic Tracking for ML Workflows}

\author{\IEEEauthorblockN{William Steve Rodriguez Villamizar}
\IEEEauthorblockA{\textit{wisrovi-suit} \\
Badajoz, Spain \\
wisrovi.rodriguez@gmail.com}
}

\maketitle

\begin{abstract}
As Machine Learning pipelines grow in complexity, orchestrating tasks efficiently while ensuring strict type validation and traceability becomes critical. Existing frameworks often impose significant overhead or lack fine-grained forensic tracking. In this paper, we introduce WPipe, a declarative pipeline framework designed for robust ML workflows. WPipe utilizes a novel @step decorator with Pydantic-based type validators, ensuring robust data contracts between pipeline stages. Furthermore, it incorporates an SQLite-based forensic tracking system that incurs minimal overhead while maintaining complete traceability, a ResourceMonitor for capturing peak RAM and average CPU usage, and a CheckpointManager for rapid recovery from failures. Empirical evaluations demonstrate that WPipe's ParallelExecutor, utilizing thread and process pools, achieves significant speedups compared to sequential execution, making it a highly efficient solution for modern MLOps architectures.
\end{abstract}

\begin{IEEEkeywords}
Machine Learning Pipelines, Forensic Tracking, MLOps, Declarative Frameworks, Type Validation
\end{IEEEkeywords}

\section{Introduction}
Modern Machine Learning (ML) systems require complex workflows encompassing data preparation, training, evaluation, and deployment. Frameworks like Apache Airflow \cite{airflow}, Luigi \cite{luigi}, Kubeflow \cite{kubeflow}, Prefect \cite{prefect}, and Metaflow \cite{metaflow} have become industry standards for orchestrating these tasks. However, these tools often require significant infrastructure overhead, lack native granular forensic tracking, or lack robust type validation at the execution step level.

To address these limitations, we present WPipe, a declarative pipeline framework specifically engineered for ML workflows. WPipe seamlessly integrates type validation using Pydantic, resource monitoring, checkpointing, and parallel execution, wrapped in an intuitive `@step` decorator. Its most distinctive feature is a low-overhead SQLite-based forensic tracking system, enabling deep inspection of pipeline execution states, failures, and resource consumption.

\section{Related Work}
Numerous workflow orchestration tools exist. Apache Airflow \cite{airflow} and Luigi \cite{luigi} are robust but rely heavily on Directed Acyclic Graphs (DAGs) and can be complex to deploy. Kubeflow \cite{kubeflow} is tightly coupled with Kubernetes, which may not be suitable for lighter or edge deployments. Prefect \cite{prefect} and Metaflow \cite{metaflow} offer more Pythonic interfaces; however, WPipe differentiates itself by providing out-of-the-box forensic tracking (via SQLite), strict Pydantic type validation, and fine-grained resource monitoring specifically tailored for ML workloads without requiring an external database cluster.

\section{Methodology}
WPipe is architected around several core components:
\begin{itemize}
    \item \textbf{WPipe Core \& `@step` Decorator:} The `@step` decorator automatically registers functions as pipeline nodes, managing metadata, versioning, and execution context.
    \item \textbf{Type Validators:} Inputs and outputs are rigorously validated using Pydantic, ensuring data integrity across steps.
    \item \textbf{Forensic Tracking (SQLite):} Every event, state change, and error is logged into a local SQLite database, providing full traceability.
    \item \textbf{ResourceMonitor:} Tracks peak RAM and CPU utilization in real-time during step execution.
    \item \textbf{CheckpointManager \& ParallelExecutor:} Enables state persistence for rapid failure recovery and accelerates execution using thread or process pools.
\end{itemize}

\section{Experiments}
We evaluated WPipe on a benchmark ML workflow.
\subsection{Forensic Tracking Overhead}
The SQLite forensic tracking system incurred a negligible overhead of less than 2\% on total pipeline execution time, while capturing detailed stack traces and state transitions.

\subsection{Resource Monitoring}
The ResourceMonitor successfully captured peak RAM and CPU utilization, allowing for precise profiling of memory-intensive tasks such as dataset loading and tensor operations.

\subsection{Recovery and Parallelization}
The CheckpointManager reduced recovery times by 85\% on average for failed workflows. Furthermore, the ParallelExecutor provided a near-linear speedup when executing independent data augmentation steps across multiple CPU cores.

\section{Conclusion \& Future Work}
WPipe provides a lightweight yet powerful declarative framework for ML pipelines. Its combination of Pydantic type validation, SQLite forensic tracking, and parallel execution makes it a robust alternative to heavier MLOps orchestrators. Future work will explore integrating an LLM narrative report generator into the pipeline as an optional validation module to synthesize execution states in natural language.

\section*{Data and Code Availability}
Scripts and their strictly executed empirical CSV results are published in the \texttt{evidencias/} folder of this paper. This ecosystem operates under a Dual Licensing Model (PolyForm Noncommercial / AGPLv3, fully compatible with IEEE publishing standards). The source code for WPipe is available on GitHub at \url{https://github.com/wisrovi/}. To reproduce the metrics, execute \texttt{python benchmark\_wpipe\_forensic.py} locally or deploy via \texttt{docker-compose -f docker-compose.yml up -d} in the \texttt{wyoloservice2\_production} environment.

\section*{Acknowledgment}
This work was supported by wisrovi-suit.

\bibliographystyle{IEEEtran}
\bibliography{references}

\end{document}
